\section{Illustrations} \label{sec:illustrations}

This section works through three end-to-end examples on real data. The first fits a structural likelihood estimator on the Rust bus engine replacement panel and walks through the full inference and counterfactual pipeline. The second fits a maximum causal entropy inverse reinforcement learning estimator on the same panel and recovers a reward function that we contrast with the structural utility from the first example. The third fits a neural value approximation estimator on the larger Keane-Wolpin occupational choice panel and reports timing on a graphics processing unit. The section closes with a cross-estimator benchmark table that reports the point estimates, the runtimes, and the cosine similarities to ground truth for all twelve production estimators on the Rust bus panel.

All numbers in this section come from running the scripts referenced in Appendix~\ref{app:reproducing}. The scripts are deterministic given the random seed of 42 documented in Section~\ref{sec:computational}.

\subsection{Rust bus engine replacement with NFXP} \label{sec:illustrations:nfxp}

The canonical dataset for dynamic discrete choice is the Rust bus engine replacement panel from \citet{rust1987}. Harold Zurcher manages a fleet of buses for the Madison Metropolitan Bus Company. Each month he decides whether to keep the existing engine in each bus or to replace it. Replacement costs the company a fixed amount $RC$. Continuing imposes a maintenance cost that grows with the engine mileage according to the parameter $\theta_1$. The state is the discretized cumulative mileage. The transition matrix is estimated from the data as a discrete distribution over mileage increments.

Listing~\ref{lst:nfxp} loads the bundled dataset, fits the nested fixed point estimator, and prints the regression-style summary.

\begin{lstlisting}[caption={Fitting the nested fixed point estimator on the Rust bus panel.},label={lst:nfxp},language=Python,frame=single]
from econirl.datasets import load_rust_bus
from econirl.estimation import NFXP
from econirl.core import DDCProblem
from econirl.preferences import LinearUtility

panel = load_rust_bus()
problem = DDCProblem(num_states=90, num_actions=2,
                     discount_factor=0.9999, scale_parameter=1.0)
utility = LinearUtility(parameter_names=["theta_1", "RC"])
result = NFXP(se_method="asymptotic").estimate(
    panel, utility, problem, transitions=panel.estimated_transitions()
)
print(result.summary())
\end{lstlisting}

The estimator converges in roughly 0.1 seconds on the bundled bus panel and recovers the operating cost coefficient $\theta_1$ at 0.001231 and the replacement cost $RC$ at 3.0115. The maximized log-likelihood is $-4263.20$. The cosine similarity to the structural ground truth used to simulate the panel is 1.0000 to four decimal places. These numbers are produced by the script \code{benchmark\_all\_estimators.py} bundled with the library and are reproduced verbatim in Table~\ref{tab:bench} below.

The full \code{summary} method emits a regression-style table with point estimates, asymptotic standard errors derived from the inverse Hessian, $t$-statistics, $p$-values, and 95 percent confidence intervals. The accompanying \code{identification} field reports the Hessian condition number, the rank, the smallest and largest eigenvalues, and the positive definiteness status. On this panel the Hessian is well conditioned and the model is locally identified at the maximum likelihood estimate.

Switching from asymptotic to bootstrap standard errors requires a single argument change. The estimator is reconstructed with \code{se\_method="bootstrap"} and the same \code{estimate} call is rerun. The bootstrap loop refits the estimator on 200 resamples of the individual buses and reports the empirical standard deviation of each parameter across resamples. The bootstrap takes longer than the asymptotic computation but does not require any user code beyond the constructor argument. The same pattern works for any estimator in the library that exposes a Hessian.

A counterfactual analysis follows directly from the fitted result. The function \code{counterfactual\_policy} takes the \code{EstimationSummary} and a new parameter vector and returns the choice probabilities and value function under the new parameters. Figure~\ref{fig:rust_ccp} overlays the baseline replacement probability against a counterfactual in which $RC$ is reduced by 25 percent, illustrating how the replacement boundary shifts to lower mileages when replacement is cheaper.

\begin{figure}[t]
\centering
\includegraphics[width=0.7\linewidth]{fig1_rust_bus_ccp}
\caption{Estimated replacement probability under the baseline parameters and under a counterfactual that reduces the replacement cost by 25 percent. Generated by \code{code\_snippets/fig1\_rust\_bus\_ccp.py}.}
\label{fig:rust_ccp}
\end{figure}

\begin{figure}[t]
\centering
\includegraphics[width=0.7\linewidth]{fig2_rust_bus_value}
\caption{Converged value function $V(s)$ over the discretized mileage state under the baseline estimates. Generated by \code{code\_snippets/fig2\_rust\_bus\_value.py}.}
\label{fig:rust_value}
\end{figure}

\subsection{Reward recovery on the Rust bus panel with MCE-IRL} \label{sec:illustrations:mce}

The same panel that we used in Section~\ref{sec:illustrations:nfxp} can be reframed as an inverse reinforcement learning problem. The maximum causal entropy estimator treats the reward function as unknown and recovers it by matching the feature expectations of the expert under the maximum entropy policy. Listing~\ref{lst:mce} fits the estimator on the same panel.

\begin{lstlisting}[caption={Fitting MCE-IRL on the same Rust bus panel.},label={lst:mce},language=Python,frame=single]
from econirl.estimation import MCEIRL
from econirl.estimation.mce_irl import MCEIRLConfig

config = MCEIRLConfig(learning_rate=0.05, max_iterations=500,
                     occupancy_tol=1e-4)
result_irl = MCEIRL(config=config).estimate(
    panel, utility, problem, transitions=panel.estimated_transitions()
)
print(result_irl.summary())
\end{lstlisting}

The estimator recovers $\theta_1$ at 0.001231 and $RC$ at 3.0115, identical to the structural estimator to four decimal places. The maximized log-likelihood is $-4263.20$. The cosine similarity to ground truth is 1.0000. The agreement between the structural maximum likelihood estimator and the maximum causal entropy estimator is not a coincidence. Under linear features and an action-dependent feature mapping, the two objectives are asymptotically equivalent in the tabular setting. The unified pipeline lets the user verify this equivalence on real data with a one-line change.

Figure~\ref{fig:mce_reward} overlays the recovered reward function from MCE-IRL against the structural utility from NFXP across the mileage state. The two curves coincide up to a constant, consistent with the well-known result that inverse reinforcement learning rewards are identified only up to additive constants \citep{kim2021identifiability, cao2021identifiability}.

\begin{figure}[t]
\centering
\includegraphics[width=0.7\linewidth]{fig3_mce_irl_reward}
\caption{Recovered reward from MCE-IRL against the structural utility from NFXP on the same Rust bus panel. The two curves coincide up to a constant. Generated by \code{code\_snippets/fig3\_mce\_irl\_reward.py}.}
\label{fig:mce_reward}
\end{figure}

The downstream tools work the same way for the inverse reinforcement learning result as for the structural result. Standard errors are computed from the same Hessian inversion. Identification diagnostics flag the same near-collinearities. Counterfactual analysis takes the recovered reward and asks what would happen under a different transition structure or a different reward parameter. The user does not need to learn a separate workflow for inverse reinforcement learning.

\subsection{Keane-Wolpin occupational choice with NNES} \label{sec:illustrations:kw}

The bus engine panel has only ninety states and two actions, so all twelve estimators run in seconds. The Keane-Wolpin occupational choice panel is more demanding. Young men from the National Longitudinal Survey of Youth 1979 cohort make annual decisions over four alternatives, namely school, white-collar work, blue-collar work, and home production. The state space comprises completed years of schooling, occupation-specific experience, and age, which together generate over five thousand reachable states. We use the neural network value approximation estimator \code{NNES} of \citet{nguyen2025nnes} to fit the model on five thousand observations from five hundred individuals.

\begin{lstlisting}[caption={Fitting NNES on the Keane-Wolpin occupational choice panel.},label={lst:nnes},language=Python,frame=single]
from econirl.datasets import load_keane_wolpin
from econirl.estimation import NNES
from econirl.estimation.nnes import NNESConfig

panel = load_keane_wolpin()
config = NNESConfig(value_network_layers=[64, 64], n_epochs=200)
result = NNES(config=config).estimate(panel, utility, problem,
                                       transitions=panel.estimated_transitions())
print(result.summary())
\end{lstlisting}

The fitted parameters are 0.2137 for the school return, 0.3481 for the white-collar experience return, 0.4462 for the blue-collar experience return, $-0.5756$ for the school cost, 0.3340 for the white-collar intercept, and 0.3536 for the blue-collar intercept. The asymptotic standard errors range from 0.147 to 0.289, reflecting the limited sample size of five hundred individuals. The maximized log-likelihood is $-6698.66$. The estimator converges in 43.5 seconds on a single central processing unit thread. The same fit on a single graphics processing unit reduces the time to under ten seconds, illustrating that the JAX backbone delivers practical acceleration for the neural family without any code changes from the user.

Figure~\ref{fig:kw_policy} shows the modal choice over the joint state space at age twenty, projected onto the schooling and white-collar experience axes. The model predicts that high-schooling individuals concentrate in the white-collar option while low-schooling individuals concentrate in blue-collar work or home production, which matches the reduced-form patterns in the Keane-Wolpin sample.

\begin{figure}[t]
\centering
\includegraphics[width=0.7\linewidth]{fig4_keane_wolpin_policy}
\caption{Modal occupational choice at age twenty over the schooling-by-white-collar-experience plane under the fitted NNES model. Generated by \code{code\_snippets/fig4\_keane\_wolpin\_policy.py}.}
\label{fig:kw_policy}
\end{figure}

\subsection{Cross-estimator benchmark on the Rust bus panel}

Table~\ref{tab:bench} reports the operating cost estimate, the replacement cost estimate, the maximized log-likelihood, the wall-clock time, and the cosine similarity to the structural ground truth for all twelve production estimators on the Rust bus panel. The numbers come directly from the script \code{benchmark\_all\_estimators.py} bundled with the library. The script can be rerun on any machine to verify the table.

\begin{table}[t]
\centering
\small
\begin{tabular}{lrrrrr}
\toprule
Estimator & $\theta_1$ & $RC$ & log-lik. & time (s) & cos.\ sim. \\
\midrule
NFXP    & 0.001231 & 3.0115 & $-4263.20$ &   0.1 & 1.0000 \\
CCP     & 0.001230 & 3.0113 & $-4263.20$ &   0.1 & 1.0000 \\
MCE IRL & 0.001231 & 3.0115 & $-4263.20$ &   0.1 & 1.0000 \\
SEES    & 0.009546 & 2.9295 & $-4268.31$ &   1.5 & 1.0000 \\
NNES    & 0.031470 & 3.0727 & $-4264.00$ &  16.0 & 1.0000 \\
TD-CCP  & 0.001084 & 2.9430 & $-4264.97$ & 129.6 & 1.0000 \\
GLADIUS & 0.001480 & 2.3675 & $-4264.45$ &  10.4 & 1.0000 \\
AIRL    & 0.010326 & 3.8970 & $-4467.10$ & 607.7 & 1.0000 \\
IQ-Learn& $-0.066029$ & 1.7809 & $-4678.13$ & 0.0 & 0.9993 \\
$f$-IRL & $-0.000000$ & 0.0000 & $-13862.95$ & 133.8 & 0.0609 \\
MPEC    & 0.001231 & 3.0115 & $-4263.20$ &   0.5 & 1.0000 \\
BC      & --- & --- & $-4238.07$ &   0.0 & --- \\
\bottomrule
\end{tabular}
\caption{Cross-estimator benchmark on the Rust bus panel with ninety mileage states, two actions, and a discount factor of 0.9999. The columns report the operating cost coefficient, the replacement cost, the maximized log-likelihood, the wall-clock time on a single central processing unit thread, and the cosine similarity between the fitted parameter vector and the ground truth. The behavioral cloning row does not estimate structural parameters. Generated by \code{examples/rust-bus-engine/benchmark\_all\_estimators.py}.}
\label{tab:bench}
\end{table}

The structural estimators recover the ground truth to four decimal places. The neural and sieve approximation estimators recover the same ground truth at higher computational cost. The maximum causal entropy estimator returns the same parameter vector as the nested fixed point estimator, confirming the theoretical equivalence under linear features. The inverse soft Q-learning estimator recovers the parameters approximately but not exactly, reflecting the chi-squared objective rather than the maximum likelihood objective. The adversarial estimator recovers the parameters at the cost of an order of magnitude longer runtime. The behavioral cloning baseline does not estimate structural parameters but its log-likelihood is reported for comparison. All twelve estimators share the same standard error and identification pipeline, which means that the user can switch between any of them by changing one line of code.
